% --- chapter
\newcommand{\chapter}[2][]{
	\newcommand{\chapname}{#2}
	\begin{flushleft}
		\begin{minipage}[t]{\linewidth}
			\includegraphics[height=1cm]{hdht-logo.png}
			\hspace{0pt}	
			\sffamily\bfseries\large Bài 35. Tính chất và cấu tạo hạt nhân
			\begin{flushleft}
				\huge\bfseries #1
			\end{flushleft}
		\end{minipage}
	\end{flushleft}
	\vspace{1cm}
	\normalfont\normalsize
}
%-----------------------------------------------------
\chapter[Độ hụt khối của hạt nhân]{Độ hụt khối của hạt nhân}
\section{Lý thuyết}
	
	Khi một hạt nhân $^A_Z X$ tạo thành từ $Z$ hạt proton và $N$ hạt nơtron thì tổng khối lượng của hạt nuclon riêng rẽ tạo thành hạt nhân là
	\begin{equation}
	m_0=Z m_p+N m_n=Z m_p+ (A-Z) m_n.
	\end{equation}
	trong đó: $m_p$ và $m_n$ lần lượt là khối lượng của các proton và nơtron.
	
	Tuy nhiên, các phép đo chính xác đã chứng tỏ rằng, khối lượng của một hạt nhân luôn nhỏ hơn tổng khối lượng của các nuclon tạo thành hạt nhân đó. Độ chênh lệch giữa hai khối lượng đó được gọi là \textbf{độ hụt khố}i của hạt nhân, kí hiệu là $\Delta m$:
	\begin{equation}
	\Delta m = m_0 - m_X = Z m_p+ (A-Z) m_n - m_X,
	\end{equation}
	trong đó:
	\begin{itemize}
		\item $m_0=Z m_p+ (A-Z) m_n$ là tổng khối lượng các nuclon lúc đầu chưa liên kết;
		\item $m_X$ khối lượng của hạt nhân $^A_Z X$;
		\item $m_p$ và $m_n$ lần lượt là khối lượng của các proton và nơtron.
	\end{itemize}

	\luuy{Khối lượng nguyên tử (ở trạng thái trung hòa):
		\begin{equation}
		m_\text{nt}=m_X+Z m_e.
		\end{equation}
	}
	
\section{Mục tiêu bài học - Ví dụ minh họa}

\begin{dang}{Tìm khối lượng của hạt nhân}

\ppgiai{
	Sử dụng công thức
	\begin{equation*}
	m_X = Z m_p+ (A-Z) m_n - \Delta m.
	\end{equation*}
}
	
	\vidu{2}{ [THPT QG 2019 - Mã đề 201] Hạt nhân $^9_4\text{Be}$ có độ hụt khối là $\SI{0,0627}{\atomicmassunit}$. Cho khối lượng của proton và nơtron lần lượt là $\SI{1,0073}{\atomicmassunit}$ và $\SI{1,0087}{\atomicmassunit}$. Khối lượng của $^9_4\text{Be}$ là
	\begin{mcq}(4)
		\item $\SI{9,0086}{\atomicmassunit}$.
		\item $\SI{9,0068}{\atomicmassunit}$.
		\item $\SI{9,0020}{\atomicmassunit}$.
		\item $\SI{9,0100}{\atomicmassunit}$.
	\end{mcq}}
	{\begin{center}
		\textbf{Hướng dẫn giải}
	\end{center}

	Khối lượng của $^9_4\text{Be}$ là
	\begin{equation*}
	m_\text{Be}= Z m_p+ (A-Z) m_n - \Delta m = 4\cdot \SI{1,0073}{\atomicmassunit} + 5 \cdot \SI{1,0087}{\atomicmassunit} - \SI{0,0627}{\atomicmassunit} = \SI{9,0100}{\atomicmassunit}.
	\end{equation*}
	
	\begin{center}
		\textbf{Câu hỏi tương tự}
	\end{center}
	
Hạt nhân $ ^7_{3} \text{Li} $ có khối lượng $ \SI{7,0144}{u} $. Khối lượng của proton và neutron lần lượt là $ \SI{1,0073}{u} $ và $ \SI{1,0087}{u} $. Độ hụt khối của hạt nhân là $ ^7_{3} \text{Li} $ là
\begin{mcq}(4)
	\item $ \SI{0,0401}{u} $.
	\item $ \SI{0,0457}{u} $.
	\item $ \SI{0,0359}{u} $.
	\item $ \SI{0,0423}{u} $.
\end{mcq}
	\textbf{Đáp án:} D.}
	
\end{dang}
	
\begin{dang}{Tìm độ hụt khối của hạt nhân.}

\ppgiai{
	Sử dụng công thức
	\begin{equation*}
	\Delta m = Z m_p+ (A-Z) m_n - m_X.
	\end{equation*}
}
	
\viduii{2}{ [THPT QG 2015 - Mã đề 138] Cho khối lượng của hạt nhân $^{107}_{\ 47}\text{Ag}$ là $\SI{106,8783}{\atomicmassunit}$; của nơtron là $\SI{1,0087}{\atomicmassunit}$; của proton là $\SI{1,0073}{\atomicmassunit}$. Độ hụt khối của hạt nhân $^{107}_{\ 47}\text{Ag}$ là 
	\begin{mcq}(4)
		\item $\SI{0,9868}{\atomicmassunit}$.
		\item $\SI{0,9868}{\atomicmassunit}$.
		\item $\SI{0,6868}{\atomicmassunit}$.
		\item $\SI{0,9686}{\atomicmassunit}$.
	\end{mcq}}
	{\begin{center}
		\textbf{Hướng dẫn giải}
	\end{center}
	
	Độ hụt khối của hạt nhân $^{107}_{\ 47}\text{Ag}$ là 
	\begin{equation*}
	\Delta m = Z m_p+ (A-Z) m_n - m_\text{Ag} = 47\cdot \SI{1,0073}{\atomicmassunit} + 60 \cdot \SI{1,0087}{\atomicmassunit} - \SI{106,8783}{\atomicmassunit} = \SI{0,9868}{\atomicmassunit}.
	\end{equation*}
	
	\begin{center}
		\textbf{Câu hỏi tương tự}
	\end{center}
	
	[THPT QG 2018 - Mã đề 209] Hạt nhân $^7_3\text{Li}$ có khối lượng $\SI{7,0144}{\atomicmassunit}$. Cho khối lượng của proton và nơtron lần lượt là $\SI{1,0073}{\atomicmassunit}$ và $\SI{1,0087}{\atomicmassunit}$. Độ hụt khối của hạt nhân $^7_3\text{Li}$ là
	\begin{mcq}(4)
		\item $\SI{0,0457}{\atomicmassunit}$.
		\item $\SI{0,0423}{\atomicmassunit}$.
		\item $\SI{0,0359}{\atomicmassunit}$.
		\item $\SI{0,0401}{\atomicmassunit}$.
	\end{mcq}

	\textbf{Đáp án:} B.}
	
\viduii{2}
{Khối lượng của nguyên tử nhôm $ \text{Al}_{13}^{27} $ là $ \SI{26,9803}{u} $. Khối lượng của nguyên tử $ H_{1}^{1} $ là $ \SI{1,007825}{u} $, khối lượng của proton là $ \SI{1,00728}{u} $ và khối lượng của neutron là $ \SI{1,00866}{u} $. Độ hụt khối của hạt nhân nhôm là
\begin{mcq}(2)
	\item $ \SI{0,242665}{u} $.
	\item $ \SI{0,23558}{u} $.
	\item $ \SI{0,023548}{u} $.
	\item $ \SI{0,023544}{u} $.
\end{mcq}
}
{
\begin{center}
	\textbf{Hướng dẫn giải}
\end{center}
Ta có:
$$
	\Delta m = 13m_{H} + 14m_{n} - m_{Al} = \SI{0,242665}{u}.
$$
\begin{center}
	\textbf{Câu hỏi tương tự}
\end{center}

Xét đồng vị Côban $ \text{Co}_{27}^{60} $ hạt nhân có khối lượng $ m_{Co} = \SI{59,934}{u} $. Biết khối lượng của các hạt $ m_{p} = \SI{1,007276}{u} $ và $ m_{n} = \SI{1,008665}{u} $. Độ hụt khối của hạt nhân đó là
\begin{mcq}(4)
	\item $ \SI{0,401}{u} $.
	\item $ \SI{0,302}{u} $.
	\item $ \SI{0,548}{u} $.
	\item $ \SI{0,544}{u} $.
\end{mcq}
\textbf{Đáp án:} C.
}
	
\end{dang}

